\documentclass[11pt,a4paper]{article}
\usepackage[utf8]{inputenc}
\usepackage[T1]{fontenc}
\usepackage[spanish,es-minimal,es-noshorthands]{babel}
\usepackage{lmodern}
\usepackage[sfdefault,lining]{FiraSans}
\usepackage[a4paper,top=2.3cm,bottom=2.2cm,left=1.8cm,right=1.8cm,headheight=1.3cm,headsep=0.4cm,footskip=1cm]{geometry}
\usepackage{amsmath,amssymb}
\usepackage{icomma}
\usepackage{graphicx}
\usepackage[table]{xcolor}
\usepackage{tikz}
\usepackage{fontawesome5}
\usepackage{ccicons}
\usepackage{fancyhdr}
\usepackage{lastpage}
\usepackage{tasks}
\usepackage{needspace}
\usepackage[most]{tcolorbox}
\usepackage{array}
\usepackage[hidelinks]{hyperref}
\AtBeginDocument{\hypersetup{pdftitle={Ficha de ejercicios de la unidad 1: Números reales. 4º ESO},pdfsubject={Matemáticas B 4º ESO},pdfkeywords={números reales, intervalos, valor absoluto, potencias, radicales, racionalización, logaritmos}}}

\definecolor{navy}{HTML}{1B2748}
\definecolor{accent}{HTML}{E8603A}
\definecolor{teal}{HTML}{1F9E90}
\definecolor{soft}{HTML}{EEF1F8}
\definecolor{muted}{HTML}{5D6685}

\setlength{\parindent}{0pt}
\setlength{\parskip}{3pt}
\settasks{label=\alph*), label-format=\bfseries\color{navy}, label-width=1.3em, item-indent=1.8em, column-sep=0.6em, after-item-skip=5pt, before-skip=3pt, after-skip=4pt}

% ---------- capçalera i peu ----------
\newcommand{\llicencia}{\href{https://creativecommons.org/licenses/by-nc-sa/4.0/deed.es}{\ccbyncsa}\ \ Licencia Creative Commons Atribución-NoComercial-CompartirIgual 4.0 Internacional (CC BY-NC-SA 4.0)}
\pagestyle{fancy}
\fancyhf{}
\fancyhead[L]{\raisebox{-0.35\height}{\includegraphics[height=0.95cm]{logo-gva.png}}}
\fancyhead[C]{\small\color{muted}Ficha de ejercicios · Unidad 1: Números reales · 4\textsuperscript{o} ESO B}
\fancyhead[R]{\raisebox{-0.35\height}{\includegraphics[height=0.95cm]{logo-vermellar.png}}}
\fancyfoot[L]{\scriptsize\color{muted}\llicencia}
\fancyfoot[R]{\small\color{muted}\thepage}
\renewcommand{\headrulewidth}{0.4pt}
\renewcommand{\headrule}{\color{muted!40}\hrule height \headrulewidth}
\fancypagestyle{primera}{\fancyhf{}\renewcommand{\headrulewidth}{0pt}
  \fancyfoot[L]{\scriptsize\color{muted}\llicencia}
  \fancyfoot[R]{\small\color{muted}\thepage}}

% ---------- nivells i exercicis ----------
\newcommand{\estrelles}[1]{{\color{accent}\ifcase#1\or\faStar\faStar[regular]\faStar[regular]\or\faStar\faStar\faStar[regular]\or\faStar\faStar\faStar\fi}}
\newcounter{exer}
\newcommand{\ex}[1]{\par\needspace{4\baselineskip}\medskip\refstepcounter{exer}\noindent{\color{navy}\bfseries\theexer.}\ \ {\small\estrelles{#1}}\ \ }

\newcommand{\seccio}[2]{\par\needspace{16\baselineskip}\bigskip
  \begin{tcolorbox}[colback=navy,colframe=navy,coltext=white,arc=2mm,boxrule=0pt,left=3mm,right=3mm,top=1.2mm,bottom=1.2mm]
  {\large\bfseries #1\quad #2}
  \end{tcolorbox}\vspace{-2pt}}

\newcommand{\exf}[1]{\par\medskip\refstepcounter{exer}\noindent{\color{navy}\bfseries\theexer.}\ \ {\small\estrelles{#1}}\ \ }
\newcolumntype{C}[1]{>{\centering\arraybackslash}p{#1}}
% quadre de recordatori (només a la fitxa d'exercicis)
\newcommand{\recorda}[1]{\par\smallskip\begin{tcolorbox}[colback=accent!7,colframe=accent!55,arc=2mm,boxrule=0.5pt,left=2.5mm,right=2.5mm,top=1mm,bottom=1mm]\small{\color{accent}\faLightbulb[regular]}\ \textbf{Recuerda:}\ #1\end{tcolorbox}}
\newcommand{\res}[1]{#1}\newcommand{\rest}[1]{#1}

\def\fmneed{11}% que el títol d'una secció no quede sol al final del full
% =====================================================================
% Capçalera FentMates de les fitxes: un color per a cada curs.
% Els colors són els de la pàgina inicial (LEVELS a assets/js/materials.js);
% el logo ve de _eines/logo/pdf/ (python3 _eines/logo/genera.py).
%
% Ús, al final del preàmbul de la unitat:
%   % =====================================================================
% Capçalera FentMates de les fitxes: un color per a cada curs.
% Els colors són els de la pàgina inicial (LEVELS a assets/js/materials.js);
% el logo ve de _eines/logo/pdf/ (python3 _eines/logo/genera.py).
%
% Ús, al final del preàmbul de la unitat:
%   % =====================================================================
% Capçalera FentMates de les fitxes: un color per a cada curs.
% Els colors són els de la pàgina inicial (LEVELS a assets/js/materials.js);
% el logo ve de _eines/logo/pdf/ (python3 _eines/logo/genera.py).
%
% Ús, al final del preàmbul de la unitat:
%   \input{../comu/fentmates.tex}
%   \fmcurs{1eso}
% i, al començament del document:
%   \capcalera{Matemàtiques}{1r ESO}{Fitxa d'exercicis · Unitat 2}{Divisibilitat}
% Opcional abans de \input: \def\fmneed{14} (espai mínim abans de cada secció).
% =====================================================================
\makeatletter
\def\fm@logodir{../../logo/pdf}
\providecommand{\fmneed}{10}
\def\fm@cursnom{}

% clau del curs, color i color del text damunt del color
\newcommand{\fm@curs}[3]{\@namedef{fm@color@#1}{#2}\@namedef{fm@text@#1}{#3}}
\fm@curs{1eso}{E6B84A}{navy}
\fm@curs{2eso}{3B82F6}{white}
\fm@curs{3eso}{FF7A50}{white}
\fm@curs{4eso-a}{A78BFA}{white}
\fm@curs{4eso-b}{6D28D9}{white}
\fm@curs{1batx-mi}{9B1C1C}{white}
\fm@curs{1batx-ccss}{EF6B6B}{white}
\fm@curs{2batx-mii}{1E6B47}{white}
\fm@curs{2batx-ccss}{2BB3A3}{white}

\newcommand{\fmcurs}[1]{%
  \@ifundefined{fm@color@#1}{\PackageError{fentmates}{Curs desconegut: #1}{}}{}%
  \edef\fm@tmp{\noexpand\definecolor{curs}{HTML}{\@nameuse{fm@color@#1}}%
               \noexpand\colorlet{curstext}{\@nameuse{fm@text@#1}}}\fm@tmp
  \colorlet{cursfosc}{curs!72!black}%
  \def\fm@logo{\fm@logodir/vertical-#1.pdf}%
  \def\fm@marca{\fm@logodir/marca-#1.pdf}}

% ---------- capçalera de la primera pàgina ----------
% #1 matèria (petit), #2 curs, #3 tipus de document i unitat, #4 títol
\newcommand{\capcalera}[4]{%
  \gdef\fm@cursnom{#2}%
  \par\noindent
  \begin{tcolorbox}[colback=curs,colframe=curs,boxrule=0pt,arc=2.2mm,width=3.5cm,height=2.15cm,
      valign=center,left=3mm,right=2mm,top=1mm,bottom=1mm,nobeforeafter,box align=center]
    \color{curstext}{\footnotesize #1}\par\vspace{1pt}{\fontsize{19}{21}\selectfont\bfseries #2\par}%
  \end{tcolorbox}\hspace{5mm}%
  \begin{minipage}[c]{\dimexpr\linewidth-3.5cm-5mm-2.9cm\relax}
    {\color{cursfosc}\small\bfseries\MakeUppercase{#3}}\par\vspace{2pt}
    {\color{navy}\fontsize{20}{23}\selectfont\bfseries #4\par}
  \end{minipage}\hfill
  \begin{minipage}[c]{2.6cm}\raggedleft\includegraphics[width=2.6cm]{\fm@logo}\end{minipage}\par
  \vspace{2pt}}

% ---------- barres de secció del color del curs ----------
\renewcommand{\seccio}[2]{\par\needspace{\fmneed\baselineskip}\bigskip
  \begin{tcolorbox}[colback=curs,colframe=curs,coltext=curstext,arc=2mm,boxrule=0pt,left=3mm,right=3mm,top=1.2mm,bottom=1.2mm]
  {\large\bfseries #1\quad #2}
  \end{tcolorbox}\vspace{-2pt}}

% ---------- estrelles de dificultat del color del curs (a Batxillerat no se'n posen) ----------
\@ifundefined{estrelles}{}{%
  \renewcommand{\estrelles}[1]{{\color{curs}\ifcase#1\or\faStar\faStar[regular]\faStar[regular]\or\faStar\faStar\faStar[regular]\or\faStar\faStar\faStar\fi}}}

% ---------- capçalera i peu de les pàgines ----------
% primera pàgina: logos institucionals al peu
\fancypagestyle{primera}{\fancyhf{}\renewcommand{\headrulewidth}{0pt}
  \fancyfoot[L]{\raisebox{-0.32\height}{\includegraphics[height=0.78cm]{logo-gva.png}}\hspace{5mm}%
    \raisebox{-0.32\height}{\includegraphics[height=0.78cm]{logo-vermellar.png}}}
  \fancyfoot[R]{\scriptsize\color{muted}\href{https://creativecommons.org/licenses/by-nc-sa/4.0/deed.ca}{\ccbyncsa}\ \ CC BY-NC-SA 4.0 · fentmates.org\qquad\small\thepage}}
% la resta: el símbol a l'esquerra i el curs a la dreta, amb el títol al centre com abans
\fancyhead[L]{\raisebox{-0.3\height}{\includegraphics[height=0.62cm]{\fm@marca}}}
\fancyhead[R]{\small\bfseries\color{cursfosc}\fm@cursnom}
\renewcommand{\headrule}{\color{curs!45}\hrule height \headrulewidth}
\makeatother

%   \fmcurs{1eso}
% i, al començament del document:
%   \capcalera{Matemàtiques}{1r ESO}{Fitxa d'exercicis · Unitat 2}{Divisibilitat}
% Opcional abans de \input: \def\fmneed{14} (espai mínim abans de cada secció).
% =====================================================================
\makeatletter
\def\fm@logodir{../../logo/pdf}
\providecommand{\fmneed}{10}
\def\fm@cursnom{}

% clau del curs, color i color del text damunt del color
\newcommand{\fm@curs}[3]{\@namedef{fm@color@#1}{#2}\@namedef{fm@text@#1}{#3}}
\fm@curs{1eso}{E6B84A}{navy}
\fm@curs{2eso}{3B82F6}{white}
\fm@curs{3eso}{FF7A50}{white}
\fm@curs{4eso-a}{A78BFA}{white}
\fm@curs{4eso-b}{6D28D9}{white}
\fm@curs{1batx-mi}{9B1C1C}{white}
\fm@curs{1batx-ccss}{EF6B6B}{white}
\fm@curs{2batx-mii}{1E6B47}{white}
\fm@curs{2batx-ccss}{2BB3A3}{white}

\newcommand{\fmcurs}[1]{%
  \@ifundefined{fm@color@#1}{\PackageError{fentmates}{Curs desconegut: #1}{}}{}%
  \edef\fm@tmp{\noexpand\definecolor{curs}{HTML}{\@nameuse{fm@color@#1}}%
               \noexpand\colorlet{curstext}{\@nameuse{fm@text@#1}}}\fm@tmp
  \colorlet{cursfosc}{curs!72!black}%
  \def\fm@logo{\fm@logodir/vertical-#1.pdf}%
  \def\fm@marca{\fm@logodir/marca-#1.pdf}}

% ---------- capçalera de la primera pàgina ----------
% #1 matèria (petit), #2 curs, #3 tipus de document i unitat, #4 títol
\newcommand{\capcalera}[4]{%
  \gdef\fm@cursnom{#2}%
  \par\noindent
  \begin{tcolorbox}[colback=curs,colframe=curs,boxrule=0pt,arc=2.2mm,width=3.5cm,height=2.15cm,
      valign=center,left=3mm,right=2mm,top=1mm,bottom=1mm,nobeforeafter,box align=center]
    \color{curstext}{\footnotesize #1}\par\vspace{1pt}{\fontsize{19}{21}\selectfont\bfseries #2\par}%
  \end{tcolorbox}\hspace{5mm}%
  \begin{minipage}[c]{\dimexpr\linewidth-3.5cm-5mm-2.9cm\relax}
    {\color{cursfosc}\small\bfseries\MakeUppercase{#3}}\par\vspace{2pt}
    {\color{navy}\fontsize{20}{23}\selectfont\bfseries #4\par}
  \end{minipage}\hfill
  \begin{minipage}[c]{2.6cm}\raggedleft\includegraphics[width=2.6cm]{\fm@logo}\end{minipage}\par
  \vspace{2pt}}

% ---------- barres de secció del color del curs ----------
\renewcommand{\seccio}[2]{\par\needspace{\fmneed\baselineskip}\bigskip
  \begin{tcolorbox}[colback=curs,colframe=curs,coltext=curstext,arc=2mm,boxrule=0pt,left=3mm,right=3mm,top=1.2mm,bottom=1.2mm]
  {\large\bfseries #1\quad #2}
  \end{tcolorbox}\vspace{-2pt}}

% ---------- estrelles de dificultat del color del curs (a Batxillerat no se'n posen) ----------
\@ifundefined{estrelles}{}{%
  \renewcommand{\estrelles}[1]{{\color{curs}\ifcase#1\or\faStar\faStar[regular]\faStar[regular]\or\faStar\faStar\faStar[regular]\or\faStar\faStar\faStar\fi}}}

% ---------- capçalera i peu de les pàgines ----------
% primera pàgina: logos institucionals al peu
\fancypagestyle{primera}{\fancyhf{}\renewcommand{\headrulewidth}{0pt}
  \fancyfoot[L]{\raisebox{-0.32\height}{\includegraphics[height=0.78cm]{logo-gva.png}}\hspace{5mm}%
    \raisebox{-0.32\height}{\includegraphics[height=0.78cm]{logo-vermellar.png}}}
  \fancyfoot[R]{\scriptsize\color{muted}\href{https://creativecommons.org/licenses/by-nc-sa/4.0/deed.ca}{\ccbyncsa}\ \ CC BY-NC-SA 4.0 · fentmates.org\qquad\small\thepage}}
% la resta: el símbol a l'esquerra i el curs a la dreta, amb el títol al centre com abans
\fancyhead[L]{\raisebox{-0.3\height}{\includegraphics[height=0.62cm]{\fm@marca}}}
\fancyhead[R]{\small\bfseries\color{cursfosc}\fm@cursnom}
\renewcommand{\headrule}{\color{curs!45}\hrule height \headrulewidth}
\makeatother

%   \fmcurs{1eso}
% i, al començament del document:
%   \capcalera{Matemàtiques}{1r ESO}{Fitxa d'exercicis · Unitat 2}{Divisibilitat}
% Opcional abans de \input: \def\fmneed{14} (espai mínim abans de cada secció).
% =====================================================================
\makeatletter
\def\fm@logodir{../../logo/pdf}
\providecommand{\fmneed}{10}
\def\fm@cursnom{}

% clau del curs, color i color del text damunt del color
\newcommand{\fm@curs}[3]{\@namedef{fm@color@#1}{#2}\@namedef{fm@text@#1}{#3}}
\fm@curs{1eso}{E6B84A}{navy}
\fm@curs{2eso}{3B82F6}{white}
\fm@curs{3eso}{FF7A50}{white}
\fm@curs{4eso-a}{A78BFA}{white}
\fm@curs{4eso-b}{6D28D9}{white}
\fm@curs{1batx-mi}{9B1C1C}{white}
\fm@curs{1batx-ccss}{EF6B6B}{white}
\fm@curs{2batx-mii}{1E6B47}{white}
\fm@curs{2batx-ccss}{2BB3A3}{white}

\newcommand{\fmcurs}[1]{%
  \@ifundefined{fm@color@#1}{\PackageError{fentmates}{Curs desconegut: #1}{}}{}%
  \edef\fm@tmp{\noexpand\definecolor{curs}{HTML}{\@nameuse{fm@color@#1}}%
               \noexpand\colorlet{curstext}{\@nameuse{fm@text@#1}}}\fm@tmp
  \colorlet{cursfosc}{curs!72!black}%
  \def\fm@logo{\fm@logodir/vertical-#1.pdf}%
  \def\fm@marca{\fm@logodir/marca-#1.pdf}}

% ---------- capçalera de la primera pàgina ----------
% #1 matèria (petit), #2 curs, #3 tipus de document i unitat, #4 títol
\newcommand{\capcalera}[4]{%
  \gdef\fm@cursnom{#2}%
  \par\noindent
  \begin{tcolorbox}[colback=curs,colframe=curs,boxrule=0pt,arc=2.2mm,width=3.5cm,height=2.15cm,
      valign=center,left=3mm,right=2mm,top=1mm,bottom=1mm,nobeforeafter,box align=center]
    \color{curstext}{\footnotesize #1}\par\vspace{1pt}{\fontsize{19}{21}\selectfont\bfseries #2\par}%
  \end{tcolorbox}\hspace{5mm}%
  \begin{minipage}[c]{\dimexpr\linewidth-3.5cm-5mm-2.9cm\relax}
    {\color{cursfosc}\small\bfseries\MakeUppercase{#3}}\par\vspace{2pt}
    {\color{navy}\fontsize{20}{23}\selectfont\bfseries #4\par}
  \end{minipage}\hfill
  \begin{minipage}[c]{2.6cm}\raggedleft\includegraphics[width=2.6cm]{\fm@logo}\end{minipage}\par
  \vspace{2pt}}

% ---------- barres de secció del color del curs ----------
\renewcommand{\seccio}[2]{\par\needspace{\fmneed\baselineskip}\bigskip
  \begin{tcolorbox}[colback=curs,colframe=curs,coltext=curstext,arc=2mm,boxrule=0pt,left=3mm,right=3mm,top=1.2mm,bottom=1.2mm]
  {\large\bfseries #1\quad #2}
  \end{tcolorbox}\vspace{-2pt}}

% ---------- estrelles de dificultat del color del curs (a Batxillerat no se'n posen) ----------
\@ifundefined{estrelles}{}{%
  \renewcommand{\estrelles}[1]{{\color{curs}\ifcase#1\or\faStar\faStar[regular]\faStar[regular]\or\faStar\faStar\faStar[regular]\or\faStar\faStar\faStar\fi}}}

% ---------- capçalera i peu de les pàgines ----------
% primera pàgina: logos institucionals al peu
\fancypagestyle{primera}{\fancyhf{}\renewcommand{\headrulewidth}{0pt}
  \fancyfoot[L]{\raisebox{-0.32\height}{\includegraphics[height=0.78cm]{logo-gva.png}}\hspace{5mm}%
    \raisebox{-0.32\height}{\includegraphics[height=0.78cm]{logo-vermellar.png}}}
  \fancyfoot[R]{\scriptsize\color{muted}\href{https://creativecommons.org/licenses/by-nc-sa/4.0/deed.ca}{\ccbyncsa}\ \ CC BY-NC-SA 4.0 · fentmates.org\qquad\small\thepage}}
% la resta: el símbol a l'esquerra i el curs a la dreta, amb el títol al centre com abans
\fancyhead[L]{\raisebox{-0.3\height}{\includegraphics[height=0.62cm]{\fm@marca}}}
\fancyhead[R]{\small\bfseries\color{cursfosc}\fm@cursnom}
\renewcommand{\headrule}{\color{curs!45}\hrule height \headrulewidth}
\makeatother

\fmcurs{4eso-b}

\begin{document}
\thispagestyle{primera}

\capcalera{Matemáticas B}{4\textsuperscript{o} ESO}{Ficha de ejercicios · Unidad 1}{\raggedright Números reales}
\vspace{10pt}
Nombre y apellidos: \rule{6.6cm}{0.4pt}\hfill Grupo: \rule{1.6cm}{0.4pt}\hfill Fecha: \rule{2.2cm}{0.4pt}
\vspace{8pt}

\begin{tcolorbox}[colback=soft,colframe=soft,arc=2mm,left=3mm,right=3mm,top=1.5mm,bottom=1.5mm]
\small Ejercicios graduados por niveles de dificultad:\quad \estrelles{1}\ básico \qquad \estrelles{2}\ medio \qquad \estrelles{3}\ avanzado.\\[2pt] Resuélvelos en el cuaderno mostrando todos los pasos. Da los resultados de manera exacta (fracciones y radicales simplificados, sin raíces en el denominador). Usa la calculadora solo cuando el enunciado lo diga. Los ejercicios 1--17 son los de la ficha del departamento; las secciones 8, 9 y 10 son de ampliación, para preparar el Bachillerato.
\end{tcolorbox}

% =====================================================
\seccio{1}{Clasificación de los números}

\exf{1} Clasifica los números siguientes en el menor conjunto numérico al que pertenecen:\begin{tasks}(4)
  \task $-\frac{15}{2}$
  \task $7,\widehat9$
  \task $\sqrt{-4}$
  \task $\pi+4$
  \task $\sqrt[3]{-8}$
  \task $\frac{16}{4}$
  \task $3,1122333\ldots$
\end{tasks}


% =====================================================
\seccio{2}{Intervalos y semirrectas}

\exf{1} Representa estos intervalos en la recta real y exprésalos como inecuaciones:\begin{tasks}(4)
  \task $A=(-\infty,\, 2]$
  \task $B=[1,\, 5]$
  \task $C=(2,\, 4]$
  \task $D=(-7,\, +\infty)$
\end{tasks}

\ex{2} Calcula las uniones e intersecciones siguientes, donde $A$, $B$, $C$ y $D$ son los intervalos del ejercicio anterior:\begin{tasks}(4)
  \task $A\cup B$
  \task $B\cup C$
  \task $A\cup D$
  \task $A\cup B\cup C$
  \task $A\cap B$
  \task $B\cap D$
  \task $A\cap C$
  \task $A\cap B\cap C$
\end{tasks}\recorda{La unión son los números que están en alguno de los intervalos; la intersección, los que están en todos a la vez.}


% =====================================================
\seccio{3}{Valor absoluto}

\exf{2} Averigua para qué valores de $x$ se cumplen las desigualdades siguientes. Expresa la solución gráficamente y con intervalos:\begin{tasks}(4)
  \task $|x|<5$
  \task $|x|\ge4$
  \task $|x+2|\le9$
  \task $|x+1|\ge6$
\end{tasks}\recorda{Si $r>0$: \ $|A|<r\iff -r<A<r$ \quad y \quad $|A|>r\iff A<-r$ o $A>r$.}


% =====================================================
\seccio{4}{Potencias}

\exf{1} Calcula las potencias siguientes:\begin{tasks}(5)
  \task $(-2)^5$
  \task $-2^5$
  \task $53^0$
  \task $(-1)^{35}$
  \task $(-3)^2$
  \task $\left(\frac12\right)^{-3}$
  \task $\left(\frac32\right)^2$
  \task $\left(\frac{-1}{2}\right)^4$
  \task $\left(\frac23\right)^{-2}$
  \task $\left(\frac15\right)^{-1}$
\end{tasks}

\ex{2} Simplifica las expresiones siguientes con potencias:\begin{tasks}(2)
  \task $\left[\left(\dfrac{-1}{3}\right)^3\cdot\left(\dfrac32\right)^{-2}\right]^2$
  \task $\dfrac{(-2)^5\cdot(2^3)^4}{-2^4\cdot(-2)^{-6}}$
  \task $\dfrac{2^4\cdot3^{-2}\cdot5^{-1}}{2^{-3}\cdot3^2\cdot5^{-2}}$
  \task $\dfrac{3^{-4}\cdot16\cdot9^{-1}}{8^2\cdot3^{-5}\cdot2^{-3}}$
\end{tasks}\recorda{Descompón las bases en factores primos y aplica $a^m\cdot a^n=a^{m+n}$, $a^m:a^n=a^{m-n}$, $(a^m)^n=a^{m\cdot n}$ y $a^{-n}=\frac1{a^n}$.}


% =====================================================
\seccio{5}{Radicales}

\exf{1} Calcula las raíces siguientes siempre que se pueda:\begin{tasks}(3)
  \task $\sqrt[3]{64}$
  \task $\sqrt[3]{-8}$
  \task $\sqrt[5]{100\,000}$
  \task $\sqrt[5]{-32}$
  \task $\sqrt{0,36}$
  \task $\sqrt{-1}$
  \task $\sqrt[3]1$
  \task $\sqrt[3]{-1}$
  \task $\sqrt1$
\end{tasks}

\ex{2} Extrae todos los factores que se pueda de las raíces siguientes:\begin{tasks}(4)
  \task $\sqrt{1280}$
  \task $\sqrt[3]{1080}$
  \task $\sqrt[3]{5184}$
  \task $\sqrt[3]{1024}$
  \task $\sqrt[4]{2^9\cdot3^{15}\cdot5^{27}}$
  \task $\sqrt{3^8\cdot5^{11}\cdot7^5}$
  \task $\sqrt{a^7\cdot b^7\cdot c^9}$
  \task $\sqrt[5]{x^{27}\cdot y^{30}\cdot z^{48}}$
\end{tasks}

\ex{2} Introduce los factores en las raíces y simplifica:\begin{tasks}(4)
  \task $\frac35\sqrt2$
  \task $2\sqrt[3]5$
  \task $2\sqrt[3]7$
  \task $4\sqrt[4]{20}$
  \task $\frac12\sqrt[4]6$
  \task $5\sqrt[3]{\frac15}$
  \task $3\sqrt[5]{15}$
  \task $\frac12\sqrt[4]{\frac12}$
\end{tasks}\recorda{El factor entra en la raíz elevado al índice: $a\sqrt[n]{b}=\sqrt[n]{a^n\cdot b}$.}

\ex{1} ¿Cuántas raíces reales hay en cada caso?\begin{tasks}(2)
  \task $\sqrt[\text{par}]{\text{positivo}}$
  \task $\sqrt[\text{par}]{\text{negativo}}$
  \task $\sqrt[\text{impar}]{\text{positivo}}$
  \task $\sqrt[\text{impar}]{\text{negativo}}$
\end{tasks}

\ex{1} Expresa en forma de una sola potencia:\begin{tasks}(4)
  \task $2^{\frac34}\cdot2^{\frac56}$
  \task $2^{\frac73}\cdot2^{-\frac52}$
  \task $\dfrac{3^{\frac12}}{3^{-\frac34}}$
  \task $\left(5^{-2}\right)^{-\frac34}$
\end{tasks}

\ex{2} Expresa con un solo radical y simplifica si es posible:\begin{tasks}(4)
  \task $\sqrt5\cdot\sqrt[3]5$
  \task $\sqrt[3]2\cdot\sqrt[4]7$
  \task $\sqrt[3]9\cdot\sqrt3$
  \task $\sqrt[4]a\cdot\sqrt[3]{a^3}\cdot\sqrt[6]{a^5}$
  \task $\dfrac{\sqrt3}{\sqrt[5]{3^4}}$
  \task $\sqrt{\sqrt[3]5}$
  \task $\sqrt[4]{\sqrt8}$
  \task $\sqrt[3]{2\sqrt{2^3}}$
\end{tasks}\recorda{Pasa a forma de potencia ($\sqrt[n]{a^m}=a^{\frac mn}$) o reduce a índice común.}

\ex{2} Calcula:\begin{tasks}(2)
  \task $\sqrt{45}-\sqrt{125}-\sqrt{20}$
  \task $\sqrt{75}-\sqrt{147}+\sqrt{675}-\sqrt{12}$
  \task $\sqrt{175}+\sqrt{63}-2\sqrt{28}$
  \task $\sqrt{20}+\frac13\sqrt{45}+2\sqrt{125}$
  \task $\sqrt{32}-\sqrt8+3\sqrt{72}-2\sqrt{200}$
  \task $7\sqrt[3]3-\sqrt[3]{24}+3\sqrt[3]{81}-4\sqrt[3]{375}$
\end{tasks}\recorda{Extrae factores y suma solo los radicales semejantes (mismo índice y mismo radicando).}


% =====================================================
\seccio{6}{Racionalización}

\exf{2} Racionaliza:\begin{tasks}(4)
  \task $\dfrac{2}{\sqrt7}$
  \task $\dfrac{1}{\sqrt3}$
  \task $\dfrac{3\sqrt2}{\sqrt3}$
  \task $\dfrac{6}{\sqrt[3]4}$
  \task $\dfrac{5}{2\sqrt[3]{10}}$
  \task $\dfrac{5}{2-\sqrt3}$
  \task $\dfrac{2\sqrt3}{\sqrt5-\sqrt3}$
  \task $\dfrac{\sqrt2}{\sqrt2+4}$
\end{tasks}\recorda{Con una suma o una resta en el denominador, multiplica por el conjugado: $(a+b)(a-b)=a^2-b^2$.}


% =====================================================
\seccio{7}{Logaritmos}

\exf{1} Calcula, utilizando la definición, los logaritmos siguientes:\begin{tasks}(4)
  \task $\log_2 8$
  \task $\log_3 81$
  \task $\log 1000$
  \task $\ln e^{-4}$
  \task $\log 0,0001$
  \task $\log_4 0,25$
  \task $\log_2\sqrt[6]4$
  \task $\ln e^3$
\end{tasks}\recorda{$\log_a b=c\iff a^c=b$. Además, $\log_a a=1$ y $\log_a 1=0$.}

\ex{1} Halla cuánto vale $x$ en cada caso:\begin{tasks}(4)
  \task $\log_3 27=x$
  \task $\log_2 x=8$
  \task $\log_x 125=3$
  \task $\ln e^5=x$
  \task $\log 0,0001=x$
  \task $\log 10\,000=x$
  \task $\log_a a=x$
  \task $\log_a 1=x$
\end{tasks}

\ex{2} Sabiendo que $\log_5 a=0,31$ y que $\log_5 b=1,4$, calcula:\begin{tasks}(2)
  \task $\log_5(a\cdot b)$
  \task $\log_5(a^3\cdot b)$
\end{tasks}


% =====================================================
\seccio{8}{Ampliación: números, intervalos, valor absoluto y potencias}

\exf{2} Escribe los números en forma de fracción y opera:\begin{tasks}(2)
  \task $2,\widehat7+3,\widehat4-1,\widehat1$
  \task $1,\widehat8+0,\widehat1$
  \task $0,1\widehat6\cdot1,\widehat5$
  \task $1,2\widehat3-0,\widehat{45}$
\end{tasks}

\ex{2} Dados $A=\left(-3,\, \frac52\right]$, $B=\left[\frac12,\, 4\right)$ y $C=(-\infty,\, 1]$, calcula:\begin{tasks}(3)
  \task $A\cap B$
  \task $A\cup C$
  \task $B\cap C$
  \task $A\cap B\cap C$
  \task $B\cup C$
\end{tasks}

\ex{2} Resuelve y representa:\begin{tasks}(4)
  \task $|2x-1|>3$
  \task $|3x+6|\le9$
  \task $|5-x|\ge2$
  \task $|x-3|<-1$
\end{tasks}

\ex{3} Escribe cada intervalo con una desigualdad con valor absoluto (busca su centro y su radio): \ $(1,\, 7)$, \ $[-5,\, 1]$, \ $(-\infty,\, -1)\cup(5,\, +\infty)$.

\ex{2} La temperatura $T$ de una nevera debe ser de $4\,^\circ$C con un error máximo de $1,5\,^\circ$C. Escríbelo con un valor absoluto y como intervalo. ¿Es aceptable una lectura de $5,7\,^\circ$C?

\ex{3} Simplifica, descomponiendo las bases en factores primos: \ $\dfrac{6^{-2}\cdot9^3\cdot2^5}{4^2\cdot27\cdot3^{-1}}$. \ Calcula también, sin calculadora: \ $\dfrac{8^{2/3}\cdot27^{-1/3}}{16^{3/4}}$, \ $\left(\dfrac1{32}\right)^{-2/5}$, \ $81^{0,75}$.


% =====================================================
\seccio{9}{Ampliación: radicales y racionalización}

\exf{2} Racionaliza el denominador de las expresiones siguientes:\begin{tasks}(4)
  \task $\dfrac{5}{\sqrt5}$
  \task $\dfrac{2\sqrt6}{\sqrt2}$
  \task $\sqrt{\dfrac52}$
  \task $\dfrac{3}{\sqrt5}$
  \task $\sqrt[3]{\dfrac32}$
  \task $\dfrac{3}{2\sqrt[3]4}$
  \task $\dfrac{3y}{x\sqrt y}$
  \task $\dfrac{x\sqrt y}{y\sqrt x}$
  \task $\dfrac{2\sqrt{12}}{5\sqrt3}$
  \task $\dfrac{3\sqrt2}{\sqrt6}$
  \task $\dfrac{3xy^2}{\sqrt[3]{x^2y}}$
  \task $\dfrac{6(3-y)}{\sqrt[3]{(3-y)^2}}$
  \task $\dfrac{-4}{\sqrt[4]3\cdot\sqrt[3]2}$
\end{tasks}

\ex{3} Racionaliza el denominador de las expresiones siguientes:\begin{tasks}(4)
  \task $\dfrac{1}{\sqrt2+\sqrt3}$
  \task $\dfrac{2}{\sqrt2-1}$
  \task $\dfrac{\sqrt2}{3-\sqrt2}$
  \task $\dfrac{1-\sqrt2}{1+\sqrt2}$
  \task $\dfrac{\sqrt5-2}{3-2\sqrt5}$
  \task $\dfrac{1}{\sqrt x-\sqrt y}$
  \task $\dfrac{\sqrt{2y}}{\sqrt2-\sqrt y}$
  \task $\dfrac{\sqrt{xy}}{2\sqrt x-3\sqrt y}$
  \task $\dfrac{5}{2(\sqrt3-\sqrt2)}$
  \task $\dfrac{\sqrt2}{2\sqrt2+1}$
  \task $\dfrac{2+\sqrt2}{2-\sqrt2}$
\end{tasks}

\ex{2} Realiza las sumas y restas siguientes:\begin{tasks}(2)
  \task $\sqrt8+\sqrt{18}+\sqrt{32}+\sqrt{50}$
  \task $\sqrt{28}+\sqrt{63}-\sqrt{175}$
  \task $6\sqrt3-4\sqrt3+5\sqrt3$
  \task $3\sqrt2+5\sqrt2-8\sqrt2$
  \task $3\sqrt2-3\sqrt8+3\sqrt{18}$
  \task $2a\sqrt3-\sqrt{27a^2}+a\sqrt{12}$
  \task $(3+a)\sqrt5-\sqrt{125}+\sqrt{5a^2}$
  \task $7\sqrt{54}-3\sqrt{18}+\sqrt{24}-\frac35\sqrt{50}-\sqrt6$
  \task $\sqrt{\frac13}+\sqrt{27}$
  \task $\sqrt{\frac12}+\sqrt2+\sqrt[4]4+\sqrt[6]8+\sqrt[4]{64}$
  \task $2\sqrt{\frac34}-\sqrt{27}+\frac14\sqrt{12}-3\sqrt{\frac{75}{9}}$
  \task $\sqrt{18}+5\sqrt{12}-2\sqrt2-\sqrt{243}$
  \task $2\sqrt3+5\sqrt[6]{27}-3\sqrt{12}$
  \task $5\sqrt{20}-2\sqrt{45}+\sqrt5$
\end{tasks}

\ex{3} Opera y simplifica:\begin{tasks}(2)
  \task $(\sqrt3+\sqrt2)^2$
  \task $(2\sqrt5-1)(2\sqrt5+1)$
  \task $\sqrt2\left(\sqrt8-\sqrt{18}\right)$
  \task $(\sqrt5+1)^2-2\sqrt5$
  \task $\dfrac{\sqrt3+1}{\sqrt3-1}-\dfrac{\sqrt3-1}{\sqrt3+1}$
  \task $\dfrac{1}{\sqrt2}+\dfrac{1}{\sqrt2+1}$
\end{tasks}

\ex{2} Un cuadrado tiene un área de $50\ \text{cm}^2$. Calcula su lado y su diagonal en forma de radical simplificado.


% =====================================================
\seccio{10}{Ampliación: logaritmos y problemas}

\exf{2} Halla $x$:\begin{tasks}(4)
  \task $\log_x 49=2$
  \task $\log_{\frac12}x=-3$
  \task $\log_x\frac18=-3$
  \task $\log_{\sqrt2}x=4$
\end{tasks}

\ex{2} Sabiendo que $\log2=0,301$ y $\log3=0,477$, calcula sin calculadora:\begin{tasks}(5)
  \task $\log12$
  \task $\log5$
  \task $\log1,5$
  \task $\log\sqrt{18}$
  \task $\log0,06$
\end{tasks}

\ex{2} Desarrolla con las propiedades de los logaritmos y escribe $E$ sin logaritmos:\begin{tasks}(2)
  \task $\log\dfrac{100x^3}{y}$
  \task $\log\sqrt{\dfrac{x}{y}}$
  \task $\log E=\log3+2\log x-\log y$
  \task $\log E=2-\log x$
\end{tasks}

\ex{1} Con la calculadora (cambio de base), calcula con tres decimales:\begin{tasks}(4)
  \task $\log_3 20$
  \task $\log_{0,5}60$
  \task $\log_2 100$
  \task $\ln7$
\end{tasks}\recorda{$\log_a x=\dfrac{\log x}{\log a}=\dfrac{\ln x}{\ln a}$.}

\ex{2} Una colonia de bacterias se duplica cada hora. ¿Cuántas horas tarda en pasar de 100 bacterias a un millón?

\ex{2} Ingresamos 1500 € en una cuenta al 4\,\% de interés compuesto anual. ¿Cuántos años deben pasar, como mínimo, para que haya 2000 €?

\ex{3} El pH de una disolución es $\text{pH}=-\log[\text{H}^+]$. Calcula el pH si $[\text{H}^+]=10^{-3}$ y si $[\text{H}^+]=2,5\cdot10^{-5}$. ¿Cuál es más ácida (pH menor)?

\ex{3} Sin calculadora, ordena de menor a mayor $\log_2 10$, \ $3$ \ y \ $\log_3 10$. Justifícalo.

\vfill
\begin{center}\color{muted}\small\faCheckCircle[regular]\ \ ¡Ánimo! Comprueba tus resultados con las herramientas interactivas de los apuntes de la unidad.\end{center}

\end{document}
